\ifdefined\gkver\else\def\gkver{student}\fi
\documentclass[\gkver]{gaokaozhenti}
\begin{document}

\chapter{2013年重庆理综}
%% paperType: 理综物理部分

\begin{enumerate}
\item
%% number: 1
%% paperName: 2013年重庆理综第 1 题 6 分
%% typeId: 1
%% type: 单选题
%% chapter: 力物体的平衡
%% point: 多力平衡
%% method: 无
%% score: 6
%% degree: 850
%% duplicateId: 0
%% body:
如图所示，某人静躺在椅子上，椅子的靠背与水平面之间有固定倾斜角 $\theta$。若此人所受重力为 $G$，则椅子各部分对他的作用力的合力大小为\xzanswer{A}

\onepicture[w=6cm, align=right]{figs/01.png}

\fourchoices[answer=A]
{$G$}
{$G\sin\theta$}
{$G\cos\theta$}
{$G\tan\theta$}

%% answer: A
%% memo:
\memoanswer{
某人静躺在椅子上受力平衡，所受合力 $F_{\text{合}}=0$，所以椅子各部分对人的作用力的合力与人的重力等大反向，故 A 选项正确。
}

\item
%% number: 2
%% paperName: 2013年重庆理综第 2 题 6 分
%% typeId: 1
%% type: 单选题
%% chapter: 原子和原子核
%% point: 核裂变
%% method: 无
%% score: 6
%% degree: 850
%% duplicateId: 0
%% body:
铀是常用的一种核燃料，若它的原子核发生了如下的裂变反应：\ce{^{235}_{92}U + ^{1}_{0}n -> a + b + 2^{1}_{0}n}，则 \ce{a + b} 可能是\xzanswer{D}

\fourchoices[answer=D]
{\ce{^{140}_{54}Xe + ^{93}_{36}Kr}}
{\ce{^{141}_{56}Ba + ^{92}_{36}Kr}}
{\ce{^{141}_{56}Ba + ^{93}_{38}Sr}}
{\ce{^{140}_{54}Xe + ^{94}_{38}Sr}}

%% answer: D
%% memo:
\memoanswer{
本题原卷及材料中未提供详解，待补充。
}

\item
%% number: 3
%% paperName: 2013年重庆理综第 3 题 6 分
%% typeId: 1
%% type: 单选题
%% chapter: 电场
%% point: 带电粒子的轨迹类问题
%% method: 无
%% score: 6
%% degree: 650
%% duplicateId: 0
%% body:
如图所示，高速运动的 $\alpha$ 粒子被位于 O 点的重原子核散射，实线表示 $\alpha$ 粒子运动的轨迹，M、N 和 Q 为轨迹上的三点，N 点离核最近，Q 点比 M 点离核更远，则\xzanswer{B}

\onepicture[w=6.5cm, align=right]{figs/03.png}

\fourchoices[answer=B]
{$\alpha$ 粒子在 M 点的速率比在 Q 点的大}
{三点中，$\alpha$ 粒子在 N 点的电势能最大}
{在重核产生的电场中，M 点的电势比 Q 点的低}
{$\alpha$ 粒子从 M 点运动到 Q 点，电场力对它做的总功为负功}

%% answer: B
%% memo:
\memoanswer{
本题原卷及材料中未提供详解，待补充。
}

\item
%% number: 4
%% paperName: 2013年重庆理综第 4 题 6 分
%% typeId: 1
%% type: 单选题
%% chapter: 运动和力
%% point: 牛顿第二定律
%% method: 无
%% score: 6
%% degree: 650
%% duplicateId: 0
%% body:
如图 1 为伽利略研究自由落体运动实验的示意图，让小球由倾角为 $\theta$ 的光滑斜面滑下，然后在不同的 $\theta$ 角条件下进行多次实验，最后推理出自由落体运动是一种匀加速直线运动。分析该实验可知，小球对斜面的压力、小球运动的加速度和重力加速度与各自最大值的比值 $y$ 随 $\theta$ 变化的图像分别对应图 2 中的\xzanswer{B}

\onepicture[w=9cm, align=right]{figs/04.png}

\fourchoices[answer=B]
{①、②和③}
{③、②和①}
{②、③和①}
{③、①和②}

%% answer: B
%% memo:
\memoanswer{
本题原卷及材料中未提供详解，待补充。
}

\item
%% number: 5
%% paperName: 2013年重庆理综第 5 题 6 分
%% typeId: 1
%% type: 单选题
%% chapter: 磁场
%% point: 带电粒子在磁场中的运动
%% method: 无
%% score: 6
%% degree: 450
%% duplicateId: 0
%% body:
如图所示，一段长方体形导电材料，左右两端面的边长都为 $a$ 和 $b$，内有带电量为 $q$ 的某种自由运动电荷。导电材料置于方向垂直于其前表面向里的匀强磁场中，内部磁感应强度大小为 $B$。当通以从左到右的稳恒电流 $I$ 时，测得导电材料上、下表面之间的电压为 $U$，且上表面的电势比下表面的电势低，由此可得该导电材料单位体积内自由运动电荷数及自由运动电荷的正负别为\xzanswer{C}

\onepicture[w=8cm, align=right]{figs/05.png}

\fourchoices[answer=C]
{$\frac{IB}{|q|aU}$，负}
{$\frac{IB}{|q|aU}$，正}
{$\frac{IB}{|q|bU}$，负}
{$\frac{IB}{|q|bU}$，正}

%% answer: C
%% memo:
\memoanswer{
本题原卷及材料中未提供详解，待补充。
}

\item
%% number: 6
%% paperName: 2013年重庆理综第 6 题 13 分
%% typeId: 4
%% type: 实验题
%% chapter: 电路
%% point: 故障分析
%% method: 无
%% score: 13
%% degree: 650
%% duplicateId: 0
%% body:
\begin{enumerate}
	\item
	我国舰载飞机在“辽宁舰”上成功着舰后，某课外活动小组对舰载飞机利用阻拦索着舰的力学问题很感兴趣。他们找来了木板、钢球、铁钉、橡皮条以及墨水，制作了如图所示的装置，准备定量研究钢球在橡皮条阻拦下前进的距离与被阻拦前速率的关系。要达到实验目的，需直接测量的物理量是钢球由静止释放时的\tkanswer[1.2cm]{高度}和在橡皮条阻拦下前进的距离，还必须增加的一种实验器材是\tkanswer[1.2cm]{刻度尺}。忽略钢球所受的摩擦力和空气阻力，重力加速度已知，根据\tkanswer[2cm]{机械能守恒（动能）}定律（定理），可得到钢球被阻拦前的速率。

	\onepicture[w=9cm, align=right]{figs/06.png}

	\item
	某同学对有故障的电热毯进行探究，题 6（2）图 \ref{2013重庆06a} 是电热毯的电路示意图，其中电热线和导线通过金属接线片连接。题 6（2）图 \ref{2013重庆06b} 为测试电路实物图，A、B 为测试表笔，电压表内阻很大，可视为理想电表。

	\threepicture[numstyle=arabic, labelA=2013重庆06a, labelB=2013重庆06b, labelC=2013重庆06c, align=right, wA=4.5cm, wB=6cm, wC=8cm]{figs/06a.png}{figs/06b.png}{figs/06c.png}

	\begin{steps}
		\item
		请在答题卡虚线框内画出与题 6（2）图 \ref{2013重庆06b} 对应的电路图。
		\item
		断开 $K_1$，用上述测试电路在 1 和 1$'$ 之间检测得知电热线无故障，然后测得电热线的 $U-I$ 曲线如题 6（2）图 \ref{2013重庆06c} 所示。已知电热线材料的电阻率为 $2.8\times10^{-7}\UOm$，电热线的直径为 $0.200\Umm$。可求得此电热线的电阻为\tkanswer[1cm]{0.6}$\UkO$，总长度为\tkanswer[1cm]{67}$\Um$。（结果均保留两位有效数字）
		\item
		为了进一步检查故障，该同学闭合开关 $K_1$ 和 $K_2$，用表笔 A 和 B 分别对题 6（2）图 \ref{2013重庆06a} 中所示的各点进行测试，部分测试结果（表中“有”“无”表示电表指针有无偏转）如下表所示。由此测试结果可判断出电路有断路，位置在\tkanswer[1.6cm]{1$'$和2$'$}（在“1 和 2”、“1$'$ 和 2$'$”、“2 和 3”、“2$'$ 和 3$'$”中选填一项）。

		\begin{center}
		\begin{tabular}{|c|c|c|c|c|c|}
		\hline
		测试点 & 3 和 3$'$ & 1 和 1$'$ & 1 和 3 & 1 和 2$'$ & 2$'$ 和 3$'$ \\ \hline
		电压表 & 有 & 有 & 无 & 有 & 无 \\ \hline
		电流表 & 无 & 有 & 有 & 无 & 有 \\ \hline
		\end{tabular}
		\end{center}
	\end{steps}
\end{enumerate}

%% answer: 
\jdanswer{
\begin{enumerate}
	\item
	高度（距水平木板的高度），刻度尺，机械能守恒（动能）

	\item
	①见解析（画出与题 6（2）图 2 对应的电路图，无答案图，待补充）；②电热线的电阻为 $0.6\UkO$，总长度为 $67\Um$；③断路位置在 1$'$ 和 2$'$。

\end{enumerate}
}
%% memo:
\memoanswer{
本题原卷及材料中未提供详解，待补充。
}

\item
%% number: 7
%% paperName: 2013年重庆理综第 7 题 15 分
%% typeId: 6
%% type: 计算题
%% chapter: 磁场
%% point: 安培力
%% method: 无
%% score: 15
%% degree: 650
%% duplicateId: 0
%% body:
小明在研究性学习中设计了一种可测量磁感应强度的实验，其装置如图所示。在该实验中，磁铁固定在水平放置的电子测力计上，此时电子测力计的读数为 $G_1$，磁铁两极之间的磁场可视为水平匀强磁场，其余区域磁场不计。直铜条 AB 的两端通过导线与一电阻连接成闭合回路，总阻值为 $R$。若让铜条水平且垂直于磁场，以恒定的速率 $v$ 在磁场中竖直向下运动，这时电子测力计的示数为 $G_2$，铜条在磁场中的长度为 $L$。
\begin{enumerate}
	\item
	判断铜条所受安培力的方向，$G_1$ 和 $G_2$ 哪个大？
	\item
	求铜条匀速运动时所受安培力的大小和磁感应强度的大小。
\end{enumerate}

\onepicture[w=6cm, align=right]{figs/07.png}

%% answer: 
\jdanswer{
\begin{enumerate}
	\item
	安培力的方向竖直向上，$G_2>G_1$

	\item
	安培力的大小 $F=G_2-G_1$，磁感应强度的大小 $B=\frac{1}{L}\sqrt{\frac{(G_2-G_1)R}{v}}$

\end{enumerate}
}
%% memo:
\memoanswer{
本题原卷及材料中未提供详解，待补充。
}

\item
%% number: 8
%% paperName: 2013年重庆理综第 8 题 16 分
%% typeId: 6
%% type: 计算题
%% chapter: 曲线运动
%% point: 圆周运动的应用
%% method: 无
%% score: 16
%% degree: 450
%% duplicateId: 0
%% body:
如图所示，半径为 $R$ 的半球形陶罐，固定在可以绕竖直轴旋转的水平转台上，转台转轴与过陶罐球心 $O$ 的对称轴 $OO'$ 重合。转台以一定角速度 $\omega$ 匀速转动，一质量为 $m$ 的小物块落入陶罐内，经过一段时间后小物块随陶罐一起转动且相对罐壁静止，它和 $O$ 点的连线与 $OO'$ 之间的夹角 $\theta$ 为 $60^{\circ}$，重力加速度大小为 $g$。
\begin{enumerate}
	\item
	若 $\omega=\omega_0$，小物块受到的摩擦力恰好为零，求 $\omega_0$；
	\item
	若 $\omega=(1\pm k)\omega_0$，且 $0<k<1$，求小物块受到的摩擦力大小和方向。
\end{enumerate}

\onepicture[w=5cm, align=right]{figs/08.png}

%% answer: 
\jdanswer{
\begin{enumerate}
	\item
	$\omega_0=\sqrt{\frac{2g}{R}}$

	\item
	当 $\omega=(1+k)\omega_0$ 时，摩擦力方向沿罐壁切线向下，大小为 $f=\frac{\sqrt{3}k(2+k)}{2}mg$；当 $\omega=(1-k)\omega_0$ 时，摩擦力方向沿罐壁切线向上，大小为 $f=\frac{\sqrt{3}k(2-k)}{2}mg$。

\end{enumerate}
}
%% memo:
\memoanswer{
本题原卷及材料中未提供详解，待补充。
}

\item
%% number: 9
%% paperName: 2013年重庆理综第 9 题 18 分
%% typeId: 6
%% type: 计算题
%% chapter: 运动和力
%% point: 动量守恒定律
%% method: 无
%% score: 18
%% degree: 450
%% duplicateId: 0
%% body:
在一种新的“子母球”表演中，让同一竖直线上的小球 A 和小球 B，从距水平地面的高度为 $ph$（$p>1$）和 $h$ 的地方同时由静止释放，如图所示。球 A 的质量为 $m$，球 B 的质量为 $3m$。设所有碰撞都是弹性碰撞，重力加速度大小为 $g$，忽略球的直径、空气阻力及碰撞时间。
\begin{enumerate}
	\item
	求球 B 第一次落地时球 A 的速度大小；
	\item
	若球 B 在第一次上升过程中就能与球 A 相碰，求 $p$ 的取值范围；
	\item
	在（2）情形下，要使球 A 第一次碰后能到达比其释放点更高的位置，求 $p$ 应满足的条件。
\end{enumerate}

\onepicture[w=3.5cm, align=right]{figs/09.png}

%% answer: 
\jdanswer{
\begin{enumerate}
	\item
	球 A 的速率 $v_0=\sqrt{2gh}$

	\item
	$1<p<5$

	\item
	$1<p<3$

\end{enumerate}
}
%% memo:
\memoanswer{
本题原卷及材料中未提供详解，待补充。
}

\item
%% number: 10
%% paperName: 2013年重庆理综第 10 题 6 分
%% typeId: 6
%% type: 计算题
%% chapter: 气体性质
%% point: 玻意耳定律
%% method: 无
%% score: 6
%% degree: 650
%% duplicateId: 0
%% body:
\begin{enumerate}
	\item
	某未密闭房间内的空气温度与室外的相同，现对该室内空气缓慢加热，当室内空气温度高于室外空气温度时\xzanswer{B}

	\fourchoices[answer=B]
	{室内空气的压强比室外的小}
	{室内空气分子的平均动能比室外的大}
	{室内空气的密度比室外的大}
	{室内空气对室外空气做了负功}

	\item
	汽车未装载货物时，某个轮胎内气体的体积为 $V_0$，压强为 $P_0$；装载货物后，该轮胎内的压强增加了 $\Delta P$。若轮胎内气体视为理想气体，其质量、温度在装载货物前后均不变，求装载货物前后此轮胎内气体体积的变化量。
\end{enumerate}

%% answer: 
\jdanswer{
\begin{enumerate}
	\item
	B

	\item
	$\Delta V=-\frac{\Delta P V_0}{P_0+\Delta P}$

\end{enumerate}
}
%% memo:
\memoanswer{
本题原卷及材料中未提供详解，待补充。
}

\item
%% number: 11
%% paperName: 2013年重庆理综第 11 题 6 分
%% typeId: 6
%% type: 计算题
%% chapter: 光学
%% point: 光的折射
%% method: 无
%% score: 6
%% degree: 650
%% duplicateId: 0
%% body:
\begin{enumerate}
	\item
	一列简谐横波沿直线传播，某时刻该列波上正好经过平衡位置的两质点相距 $6\Um$，且这两质点间的波峰只有一个，则该简谐横波可能的波长为\xzanswer{C}

	\fourchoices[answer=C]
	{$4\Um$、$6\Um$ 和 $8\Um$}
	{$6\Um$、$8\Um$ 和 $12\Um$}
	{$4\Um$、$6\Um$ 和 $12\Um$}
	{$4\Um$、$8\Um$ 和 $12\Um$}

	\item
	利用半球柱形玻璃，可减小激光光束的发散程度。在如图所示的光路中，A 为激光的出射点，O 为半球柱形玻璃横截面的圆心，AO 过半圆顶点。若某条从 A 点发出的与 AO 成 $\alpha$ 角的光线，以入射角 $i$ 入射到半圆弧上，出射光线平行于 AO，求此玻璃的折射率。

	\onepicture[w=6cm, align=right]{figs/11.png}
\end{enumerate}

%% answer: 
\jdanswer{
\begin{enumerate}
	\item
	C

	\item
	$n=\frac{\sin i}{\sin(i-\alpha)}$

\end{enumerate}
}
%% memo:
\memoanswer{
本题原卷及材料中未提供详解，待补充。
}

\end{enumerate}

\end{document}
